\chapter*{Conclusion générale}
\addcontentsline{toc}{chapter}{Conclusion générale}

Ce projet de fin d'études a porté sur la conception et la mise en œuvre d'une solution décisionnelle complète au service du CETUD, visant à exploiter l'Enquête Ménages-Déplacements et les données de comptage de trafic de Dakar pour éclairer la prise de décision en matière de mobilité urbaine.

La démarche suivie s'est appuyée sur la méthodologie GIMSI, qui a structuré le travail en quatre grandes étapes~: l'identification des décideurs et de leurs besoins, la conception d'un modèle décisionnel en constellation, la mise en œuvre technique de la chaîne ETL avec Talend et de l'entrepôt de données sous PostgreSQL, puis la restitution des indicateurs à travers le cockpit Knowage BI.

Au-delà de ce volet purement décisionnel, le projet a été enrichi par le développement de quatre modèles de Machine Learning~: la segmentation des usagers par K-Means, la recommandation de mode de transport par Random Forest, la détection d'anomalies de trafic par un consensus de trois algorithmes (Z-Score, Isolation Forest, Local Outlier Factor) complété par un clustering spatial DBSCAN, et la prédiction du risque d'inaccessibilité par Gradient Boosting. Ces modèles ont été exposés au sein d'une application web complète, combinant un backend FastAPI et un frontend React, à destination à la fois des citoyens et des décideurs du CETUD.

Sur le plan des compétences acquises, ce projet a permis une mise en pratique transversale couvrant l'ensemble de la chaîne de valeur de la donnée~: de la modélisation décisionnelle classique (schéma en constellation, ETL, entrepôt de données) jusqu'aux techniques avancées de Machine Learning (apprentissage supervisé et non supervisé, détection d'anomalies), en passant par le développement d'une application web complète. Cette double dimension, décisionnelle et prédictive, constitue l'apport principal du projet par rapport à une solution de reporting classique.

Le système présente toutefois certaines limites, assumées au regard du cadre d'un projet de fin d'études~: l'absence de données temps réel, l'utilisation de l'EMD 2019 (la plus récente disponible publiquement), une authentification simplifiée côté client, et l'absence de tests automatisés. Ces limites ouvrent autant de pistes d'amélioration pour une mise en production~: intégration de données temps réel issues des opérateurs de transport, conteneurisation Docker, authentification renforcée (JWT, bcrypt), mise en place de tests automatisés et de pipelines d'intégration continue, et, à plus long terme, exploration de techniques de Deep Learning (réseaux de neurones récurrents pour la prédiction de trafic, réseaux de neurones sur graphes pour la modélisation du réseau de transport).

En définitive, ce projet démontre la faisabilité et l'intérêt d'une approche combinant Business Intelligence traditionnelle et Machine Learning pour répondre aux enjeux de mobilité urbaine d'une métropole comme Dakar, et pose les bases d'un système qui pourra être consolidé et étendu dans le cadre d'une collaboration future avec le CETUD.

\clearpage
